\documentclass[11pt]{article}
\usepackage{ap-lang}
\title{Part 4: Style}
\author{Paul Wang}
\date{November 2026}
\docpart{Part 4 \textperiodcentered\ Study Guide}

\begin{document}
\maketitle

\begin{abstract}
Style is how a writer says what they say: the words, the shape of the sentences, the comparisons, and the attitude they convey. On this exam style is never studied for its own sake; every question asks how a stylistic choice serves the argument or affects the audience. This guide covers diction (Section 1), syntax (Section 2), comparisons and figurative language (Section 3), tone, shifts, and irony (Section 4), punctuation and concision in writing questions (Section 5), and how style earns credit in the essays (Section 6). Each section ends with check questions; the key is on the last page.
\end{abstract}

\scoringline{Check questions: \blank\ / \qtotal \qquad Target: 13+ \qquad Terms I mix up: \blank[10em]}

\section{Diction}

\textbf{Diction} is word choice. Every word has a \textbf{denotation} (dictionary meaning) and a \textbf{connotation} (associations and feeling). ``Thrifty,'' ``frugal,'' and ``stingy'' denote the same habit; their connotations range from praise to insult.

\begin{terms}
Register & The level of formality: formal, conversational, colloquial, technical. A shift in register is often a deliberate signal. \\
Concrete vs.\ abstract & ``The cracked basin'' vs.\ ``infrastructure problems.'' Concrete diction makes an argument vivid and believable. \\
Specific vs.\ general & ``Eleven children by the feed store'' vs.\ ``many children.'' Specificity builds credibility. \\
Loaded language & Words chosen for strong connotations (``scheme,'' ``slaughter,'' ``hero''). \\
Euphemism & A mild word for a harsh fact (``downsizing'' for firing). Writers point out euphemisms to expose what they hide. \\
\end{terms}

\begin{trap}
On word-choice questions, the answer is about the \emph{effect} of the connotation in this context, not the dictionary meaning. Ask: why this word and not its neutral synonym?
\end{trap}

\begin{q}{3}{4.1 Diction}{``Scheme'' carries a negative connotation of deceit that ``plan'' lacks.}
A writer refers to a city's redevelopment proposal as a ``scheme'' rather than a ``plan.'' The word choice primarily conveys
\ch{confusion about the proposal's details}{a neutral description of the proposal}{suspicion that the proposal is deceptive or self-serving}{approval of the proposal's ambition}
\end{q}

\begin{q}{4}{4.1 Diction}{Quoting the euphemism exposes the harsh fact it softens.}
A writer puts ``workforce optimization'' in quotation marks and follows it with ``which meant that 340 people lost their jobs on a Friday.'' The writer's purpose is mainly to
\ch{show the writer's expertise in business}{define a technical term for a general audience}{praise the company's efficiency}{expose a euphemism by stating the harsh fact it conceals}
\end{q}

\section{Syntax}

\textbf{Syntax} is the arrangement of words into sentences: length, order, and structure.

\begin{terms}
Sentence variety & A short sentence after long ones lands with force. ``It did not.'' \\
Parallelism & Similar grammatical structures in a series: ``to teach, to heal, to build.'' Creates rhythm and equal weight. \\
Antithesis & Contrasting ideas in parallel form: ``not a punishment but a loan.'' \\
Anaphora & Repeating words at the start of successive clauses: ``It buys\ldots It buys\ldots It buys\ldots.'' \\
Periodic sentence & Delays the main clause to the end, building suspense: ``Through floods, through fires, through two wars, the bridge stood.'' \\
Cumulative (loose) sentence & Main clause first, then added details: ``The bridge stood, through floods, through fires, through two wars.'' \\
Asyndeton / polysyndeton & Omitting conjunctions (``We came, we saw, we left'') for speed; adding them (``and the rain and the mud and the cold'') for weight. \\
Rhetorical question & A question asked for effect, not an answer; invites the reader to supply the writer's conclusion. \\
Fragment & An intentional incomplete sentence for emphasis: ``Not once.'' \\
Juxtaposition & Placing two things side by side so the contrast speaks. \\
\end{terms}

\begin{keyidea}[How to answer syntax questions]
Name the structure, then its effect on the meaning: the repetition \emph{emphasizes} X; the short sentence \emph{punctuates} Y; the parallel list \emph{gives equal weight} to Z. The device alone is never the answer.
\end{keyidea}

\begin{q}{4}{4.2 Syntax}{Contrasting ideas in balanced, parallel form = antithesis.}
``We do not inherit the land from our parents; we borrow it from our children'' is an example of
\chh{an analogy}{asyndeton}{a periodic sentence}{antithesis}
\end{q}

\begin{q}{2}{4.2 Syntax}{A short sentence after long ones lands with emphasis.}
After three long sentences describing a family's careful preparations for a flood, a writer adds: ``It was not enough.'' The short sentence primarily serves to
\ch{introduce a new topic}{emphasize, by contrast with the long sentences before it, how futile the preparations were}{suggest the writer is uncertain}{provide statistical evidence}
\end{q}

\begin{q}{1}{4.2 Syntax}{Holding the main clause to the end builds suspense.}
``Through two wars, a depression, and the closing of every mill in the valley, the church kept its doors open.'' This is best described as
\chh{a periodic sentence}{a cumulative sentence}{a rhetorical question}{a fragment}
\end{q}

\section{Comparisons and figurative language}

\begin{terms}
Metaphor / simile & An implied or explicit comparison. Ask what qualities transfer from one thing to the other. \\
Analogy & An extended comparison used to explain or argue: the economy as a household budget. Watch for where the analogy breaks down. \\
Allusion & A reference to a well-known person, event, or text (the Bible, Shakespeare, history). It borrows associations and builds shared ground with an audience that recognizes it. \\
Imagery & Sensory detail that lets the reader see, hear, or feel. \\
Personification & Giving human qualities to the nonhuman (``the grass has been trying to tell us''). \\
Hyperbole / understatement & Deliberate exaggeration or deliberate minimizing, often for humor or irony. \\
\end{terms}

\begin{trap}
Allusions on this exam assume knowledge an educated audience shares, but questions ask about their \emph{function} (what associations they bring), not about the original source in detail.
\end{trap}

\begin{q}{2}{4.3 Comparisons}{The metaphor transfers a ticking clock's urgency to the issue.}
A writer calls a town's aging water system ``a clock that is already ticking.'' The metaphor primarily conveys
\ch{that the system works with precision}{that failure is approaching and delay is costly}{that the system is old-fashioned and charming}{that the town should buy new clocks}
\end{q}

\begin{q}{4}{4.3 Allusion}{An allusion borrows the associations of the story it recalls.}
A senator calls a lone holdout on a committee ``our David among the Goliaths.'' The allusion primarily serves to
\ch{criticize the holdout for stubbornness}{establish the senator's religious beliefs}{suggest that the holdout will lose}{cast the holdout as an underdog with the moral advantage}
\end{q}

\section{Tone, shifts, and irony}

\textbf{Tone} is the writer's attitude toward the subject or audience, conveyed through diction, syntax, and detail. Name tone precisely and in pairs when it is mixed.

\begin{compare}{Family}{Useful tone words}{Signals}
Approving & admiring, reverent, affectionate, celebratory, earnest & praise, warm detail, ``we'' \\
Critical & skeptical, wry, sardonic, indignant, disapproving & quotation marks around others' terms, understatement, irony \\
Measured & reflective, candid, pragmatic, judicious, conciliatory & concessions, qualifiers, balanced syntax \\
Urgent & insistent, impassioned, exhortative, grave & imperatives, short sentences, repetition \\
\end{compare}

A \textbf{shift} in tone is a turning point: from humor to gravity, from criticism to hope, from personal to public. Questions ask what the shift signals about the argument.

\textbf{Irony} is a gap between what is said and what is meant (\textbf{verbal irony}), or between what is expected and what happens (\textbf{situational irony}). \textbf{Satire} uses irony and exaggeration to criticize.

\begin{q}{3}{4.4 Tone}{Mock praise that means the opposite = verbal irony.}
After describing a highway project that took nine years and doubled its budget, a writer remarks, ``a triumph of planning.'' The remark is an example of
\chh{a euphemism}{hyperbole}{verbal irony}{an allusion}
\end{q}

\begin{q}{2}{4.4 Tone}{A shift from humor to gravity signals the turn to the serious point.}
A speech opens with jokes about the town's terrible weather, then turns: ``But last winter, the weather was not funny.'' The shift in tone primarily signals
\ch{that the audience disliked the jokes}{the move from the speaker's rapport with the audience to the serious issue she wants to address}{a change of speaker}{that the speaker has run out of material}
\end{q}

\section{Punctuation, concision, and style in writing questions}

\subsection{Punctuation for effect}
\begin{terms}
Colon & Introduces what fulfills the preceding clause: an explanation, list, or punchline. ``There was one problem: no one had asked the tenants.'' \\
Dash & Sets off an interruption or a sharp turn with emphasis; a pair of dashes isolates a phrase more forcefully than commas. \\
Semicolon & Joins two closely related independent clauses, often balancing them. \\
Parentheses & De-emphasize: an aside, often wry or confiding. \\
\end{terms}

\subsection{Writing-question style}
Style questions on the writing sets ask you to choose the version that is \textbf{precise} (the exact word for the meaning), \textbf{concise} (no redundancy: not ``past history'' or ``each and every''), \textbf{consistent in register} (no slang in a formal letter), and \textbf{fitted to the stated goal} (``to convey the writer's frustration,'' ``to maintain a formal tone'').

\begin{q}{1}{4.5 Writing: style}{The concise version keeps the meaning without redundancy.}
Which version is the most concise without losing meaning?
\ch{The council approved the budget on Tuesday.}{The budget was, on Tuesday, approved by the council, which approved it.}{The council, on the day of Tuesday, gave its approval to the budget.}{On Tuesday, the council approved and passed the budget, giving it approval.}
\end{q}

\begin{q}{1}{4.5 Writing: style}{Slang breaks the formal register of a letter to officials.}
In a formal letter to the state transportation board, a student writes: ``The bus schedule is honestly kind of a mess.'' Which revision best maintains a formal tone?
\ch{The bus schedule is inconsistent and difficult for riders to follow.}{The bus schedule is a total disaster, no cap.}{Honestly, the bus schedule is a mess.}{The bus schedule is kind of messy, honestly.}
\end{q}

\section{Style in the essays}

In the \textbf{rhetorical analysis} essay, analyze style as a set of choices with effects: not ``she uses imagery'' but ``the image of the empty swing set makes the cost of closing the daycare visible to council members who have only seen it as a line item.''

In \textbf{all three} essays, your own style can contribute to Row C when it is \emph{consistently} vivid and persuasive: precise verbs, varied syntax, a controlled tone. One dramatic sentence is not enough; the point is earned across the essay. Plain, clear, precise prose always beats ornate prose that blurs the argument.

\begin{q}{3}{4.6 Style in essays}{Analysis links the stylistic choice to its effect and purpose.}
Which sentence from a rhetorical analysis essay best analyzes style?
\ch{``The writer's style is very effective and persuasive.''}{``The passage has many long sentences and some short ones.''}{``By ending each paragraph with a four-word sentence, the writer makes every argument land like a verdict, leaving the board nothing to debate.''}{``The writer uses a lot of imagery and metaphors.''}
\end{q}

\makekey
\end{document}
